\section{Certified numerical enclosure}\label{sec:enclosure}
\begin{theorem}[Strict rational bracket]\label{thm:O8}
For the rational numbers \(\ell,h\) displayed after Theorem~\ref{thm:main},
\[
 \ell<\St<h,
\]
with the exact width in \eqref{eq:width}.
\end{theorem}
\begin{proof}
The degree-48 analytic enclosure of Certificate~\ref{cert:local}
represents the same germ as the degree-36 global proof. It is used only
after the global equivalence and uniqueness theorems. Put
\(f(s)=\Res(a(s))\); then
\[
 f'(s)=\frac{\Res'(a(s))}{\sigma'(a(s))}>0.
\]
The exact directed interval for this quotient is recorded in
\pkg{scalar.json}. Starting at the full closed interval \([L,U]\),
the certificate performs 13 inclusions
\begin{equation}\label{eq:newton}
 J_{k+1}\supseteq
 J_k\cap\left(m_k-\frac{[f](m_k)}{[f']([L,U])}\right),\qquad
 m_k\in J_k.
\end{equation}
Every recorded output lies inside its input and strictly shrinks it.
The mean value theorem shows that the root remains in each exact
intersection; rational verification checks that the stored rounded
output contains that intersection. The last output is inflated by
exactly \(10^{-52}\) on each side to give \([\ell,h]\).

The endpoint residual enclosures, generated for \(f\) at these
\emph{\(s\)-valued} endpoints, have the following outward rational
summaries (the full decimal intervals are in \pkg{root.json}):
\begin{align*}
 [f](\ell)&\subseteq[-5.720120,-5.720084]\times10^{-51},\\
 [f](h)&\subseteq[5.804107,5.804143]\times10^{-51}.
\end{align*}
Their strict signs and continuity place a zero strictly between
\(\ell,h\). Theorem~\ref{thm:O6} and scalar equivalence identify it with
the unique global matching parameter; Theorem~\ref{thm:O7} identifies
that parameter with \(\St\). Exact integer subtraction proves
\eqref{eq:width}, and rational comparisons prove \(L<\ell<h<U\).
\end{proof}

The finite check of \eqref{eq:newton} uses integer/Fraction arithmetic,
independently of the interval code's midpoint helper: each recorded center
is verified as an exact rational point inside its own input. The
residual and derivative enclosures still depend on the analytic-germ
proof and the interval arithmetic. Checking the final rational arithmetic
alone would not establish those analytic enclosures.

Scouting suggested computational centers and a central phase interval.
Proof acceptance instead uses rational identities, directed interval
bounds, full covers and strict signs. Historical root digits, stored
success flags and agreement of floating-point optimizers supply no
premise of this enclosure.
